% GENERATED by diagnostics/build_rc_tables.R - do not hand-edit.
\begin{table}[htbp]\centering
\begin{threeparttable}
\caption{Is the rise a product of reports becoming more quantified?}
\begin{tabular}{lcc}
\toprule
 & Year term & Period terms \\
\midrule
Year (per decade), unadjusted & +0.439$^{***}$ & \\
Year (per decade), adjusted & +0.284$^{*}$ & \\
\quad \emph{attenuation} & \emph{35\%} & \\
\addlinespace
1997--2004 (vs 1984--96) & & +1.13$^{**}$ \\
2005--13 (vs 1984--96) & & +0.81 \\
2014--25 (vs 1984--96) & & +1.03$^{*}$ \\
\addlinespace
Digit density & \multicolumn{2}{c}{+46.2$^{***}$} \\
\bottomrule
\end{tabular}
\begin{tablenotes}[flushleft]\footnotesize
\item Logistic models of whether a report is coded a serious delivery failure. $n=416$ reports,
stratified by period. Adjusted models control for digit density, percentage and over-run language,
verdict vocabulary, and whether the extract carried an explicit verdict. Coefficients are log-odds.
\item \textbf{The mechanism is real and the trend survives it.} Digit density strongly predicts a
serious code, so reports containing more numbers are likelier to be recorded as serious. Adjusting
for it removes about 35 per cent of the year effect and leaves the remainder.
\item \textbf{The adjusted shape differs from the unadjusted shape.} With period terms, all three
post-1996 periods sit above 1984--96 but are not distinguishable from one another: adjusted, the
pattern reads as a step change after the mid-1990s rather than a compounding rise.
\item Currency mentions were measured and discarded: the pound sign does not survive extraction
from older and image-scanned reports, so the measure was zero for every pre-1997 report.
\item $^{*}p<.05$; $^{**}p<.01$; $^{***}p<.001$.
\item Source: \texttt{pdf\_era/text\_features\_sample.py}, models in
\texttt{diagnostics/fig214\_supporting\_tests.R}.
\end{tablenotes}
\end{threeparttable}
\label{tab:ch02naoquant}
\end{table}
